% from ega4-8-fr.tex
\oldpage[IV-3]{5}
\phantomsection
\label{chapter:IV3-fr}

\begin{center}
{\large CHAPITRE IV (suite)\par}
\vspace{1.5em}
{\Large\bfseries ÉTUDE LOCALE DES SCHÉMAS\\
ET DES MORPHISMES DE SCHÉMAS\par}
\end{center}

\setcounter{section}{7}
\section{Limites projectives de préschémas}
\label{section:IV.8-fr}

\setcounter{subsection}{0}
\subsection{Introduction}
\label{subsection:IV.8.1-fr}

\begin{env}[8.1.1]
\label{IV.8.1.1-fr}
Dans ce paragraphe, nous allons étudier systématiquement la situation
suivante. Soient $I$ un ensemble préordonné filtrant croissant,
$(A_\lambda,\varphi_{\mu\lambda})$ un système inductif d'anneaux ayant $I$
pour ensemble d'indices, $A=\varinjlim A_\lambda$ sa limite inductive. Pour
tout $\alpha\in I$ et tout $A_\alpha$-préschéma $X_\alpha$, considérons les
$A_\lambda$-préschémas
$X_\lambda=X_\alpha\otimes_{A_\alpha}A_\lambda$ pour $\lambda\geq\alpha$,
et le $A$-préschéma $X=X_\alpha\otimes_{A_\alpha}A$~; il est clair que les
préschémas $X_\lambda$ (pour $\lambda\geq\alpha$) forment un \emph{système
projectif}, et on verra (8.2.5) que $X$ est une \emph{limite projective} de ce
système dans la catégorie des préschémas. Nous nous proposons de trouver des
conditions sur $X_\alpha$ ou sur les $A_\lambda$ permettant de démontrer des
propriétés du type suivant~: pour que $X$ possède une propriété $P$ (par
exemple, la propriété d'être propre sur $S=\operatorname{Spec}(A)$, ou
irréductible, ou connexe, etc.), il faut et il suffit qu'il existe un indice
$\mu\geq\alpha$ tel que, pour tout $\lambda\geq\mu$, $X_\lambda$ ait (par
rapport à $S_\lambda=\operatorname{Spec}(A_\lambda)$, le cas échéant), la
même propriété $P$. Nous obtiendrons des énoncés analogues pour des propriétés
de $\mathcal O_X$-Modules, de $A$-morphismes de $A$-préschémas, etc. Nous
montrerons également (8.9.1) que la donnée d'un $A$-préschéma de présentation
finie (1.6.1) équivaut essentiellement à la donnée d'un $A_\lambda$-préschéma
de présentation finie $X_\lambda$ pour un $\lambda$ assez grand, $X$ étant
alors \emph{isomorphe à} $X_\lambda\otimes_{A_\lambda}A$. On a aussi des
énoncés analogues pour les $\mathcal O_X$-Modules de présentation finie,
leurs homomorphismes, les $A$-morphismes de $A$-préschémas de présentation
finie, etc.
\end{env}

\begin{env}[8.1.2]
\EGAReaderTarget{ega4.c04.statement.8-1-2}{IV.8.1.2-fr}
L'utilité de tels résultats apparaîtra par exemple dans les questions
suivantes~:
\begin{enumerate}
\item[\emph{a)}] Soient $Y$ un préschéma, $y$ un point de $Y$, $(U_\lambda)$
le système projectif filtrant décroissant des voisinages ouverts affines de
$y$ dans $Y$~; si $A_\lambda$ est l'anneau de $U_\lambda$, les $A_\lambda$
forment un système inductif filtrant croissant, dont la limite inductive $A$
est l'anneau local $\mathcal O_y$~; d'ailleurs, si on désigne par
$\mathfrak p_\lambda$ l'idéal premier $\mathfrak j_y$ dans l'anneau
$A_\lambda$, le système inductif $(A_\lambda)$ est cofinal à tout système
inductif $(A_\alpha)_f$, où $f$ parcourt
$A_\alpha-\mathfrak p_\alpha$ (pour un $\alpha$ \emph{fixé}), car les $D(f)$
forment un système fondamental de voisinages de $y$ dans $U_\alpha$,
